\section{ A Unified Classification}
The various situations may be summarized as follows.
\begin{center}
\begin{tabular}{p{0.25\textwidth}|p{0.3\textwidth}|p{0.3\textwidth}}
\hline
Mechanism & What remains symmetric? & What loses symmetry?\
\hline
Explicit breaking &
Possibly nothing at the level of the broken group &
Action and/or equations\
\hline
Boundary-induced reduction &
Local action and equations &
Physical boundary-value problem\
\hline
Initial-data reduction &
Equations and solution space &
Chosen solution and initial-value problem\
\hline
Ordinary asymmetric solution &
Equations and solution space &
Individual solution\
\hline
Spontaneous symmetry breaking &
Action and equations &
Vacuum or equilibrium state\
\hline
Gauge fixing &
Gauge-invariant theory &
Particular mathematical representation\
\hline
Symmetry enhancement &
Limiting theory &
Nearby finite-parameter theories have smaller symmetry\
\hline
\end{tabular}
\end{center}
This table makes clear why the phrase ``symmetry is broken'' is incomplete unless its precise mechanism is specified.

